%% ============================================================
%%  Section: Conclusion
%%  Depositor Discipline in Modern Banking (2026-04-04 draft)
%%  Target length: ~1 page
%%
%%  New bib entry required:
%%    correia2026bank  (see suggested entry below)
%% ============================================================

\section{Conclusion}
\label{sec:conclusion}

This paper provides causal evidence that uninsured depositors discipline
community banks in response to credit-risk disclosures in a non-crisis setting.
Exploiting the CECL Day-One retained-earnings adjustment as a predetermined
information shock, we show that community banks in the top decile of CECL
exposure experienced a 39-basis-point contraction in uninsured time deposits,
a simultaneous 71-basis-point increase in insured time deposits, and a
125-basis-point widening of the within-bank insured-uninsured deposit composition
gap---all relative to low-CECL peers after the 2023Q1 adoption date.  A
triple-difference specification with bank $\times$ quarter fixed effects confirms
that this is a composition shift within banks rather than a uniform funding
retreat, ruling out aggregate funding shocks as an alternative explanation.
The deposit reallocation carried a real financial cost: high-CECL banks paid
8--9 basis points more in annualized funding costs, experienced a 15-basis-point
drag on return on assets, and saw net interest margins compress by approximately
9 basis points.  Heterogeneity by depositor sophistication reveals a two-channel
mechanism: sophisticated depositors extract rate premiums immediately upon
disclosure and reduce balances gradually as CDs mature, amplifying the financial
penalties at institutions whose depositor bases are most capable of acting on
credit-risk information.

These findings have direct implications for ongoing policy debates about deposit
insurance reform.  The evidence establishes that uninsured depositor discipline
remains an active governance force in U.S. community banking---even after a
decade of post-crisis rescue interventions, emergency liquidity facilities, and
repeated episodes in which the FDIC extended de facto coverage to uninsured
creditors at failing institutions.  The discipline mechanism has not been
extinguished; it operates whenever depositors receive credible, unambiguous
information about bank credit risk.  This matters for proposals to substantially
raise the \$250{,}000 FDIC coverage threshold or to introduce broad-based
unlimited deposit insurance for business accounts.  By moving a larger share of
deposits inside the insurance perimeter, such reforms would directly attenuate
the monitoring and pricing pressures we document, weakening the market
discipline that complements regulatory oversight as a check on risk-taking
\citep{nier2006market, demirgucc2004market, flannery2019market}.

The welfare case for higher limits is not straightforward.  Proponents argue
that broader coverage reduces the risk of self-fulfilling runs, lowering
systemic fragility.  This argument has force in models where solvent banks can
be brought down by coordination failures alone.  Recent historical evidence,
however, suggests that the link between runs and failures is considerably more
conditional.  \citet{correia2026bank}, examining 3,421 bank runs in the United
States between 1863 and 1934, find that failures following runs were concentrated
almost exclusively at institutions with poor fundamentals; strong banks survived
runs through interbank cooperation and liquidity support.  If runs are
disciplinary rather than purely destabilizing---occurring disproportionately at
weak banks and leading to failure only when the bank was already in trouble---then
eliminating the threat of runs by extending deposit insurance primarily removes
a governance mechanism from the banks that need it most, without commensurate
benefits in systemic stability.  Taken together, our results and this broader
evidence suggest that substantially raising deposit insurance limits carries
meaningful moral-hazard costs, while the stability benefits may be smaller than
commonly assumed.

%% ============================================================
%%  SUGGESTED NEW BIB ENTRY (add to depositor-discipline-references.bib)
%%
%%  @techreport{correia2026bank,
%%    title   = {Bank Runs With and Without Bank Failure},
%%    author  = {Correia, Sergio and Luck, Stephan and Verner, Emil},
%%    year    = {2026},
%%    note    = {arXiv preprint arXiv:2601.20285}
%%  }
%%
%% ============================================================
